
\section{Artifact Grammar}


\begin{figure}[H]
\centering
\begin{tikzpicture}[every node/.style={figlabel}, box/.style={figbox, minimum width=2.85cm, minimum height=.92cm, align=center}]
\path[figpanel] (-1.25,-3.45) rectangle (8.65,3.05);
\node[figtitle, anchor=west] at (-0.8,2.55) {ARTIFACT CLASSES};
\node[box,fill=BitBlue!10,draw=BitBlue!65] (boot) at (0,1.55) {Boot};
\node[box,fill=BitTeal!10,draw=BitTeal!70] (state) at (3.7,1.55) {State};
\node[box,fill=BitGold!14,draw=BitGold!76!black] (task) at (7.4,1.55) {Task};
\node[box,fill=BitRose!10,draw=BitRose!62] (hist) at (0,-0.35) {History};
\node[box,fill=BitBlue!7,draw=BitBlue!62] (mem) at (3.7,-0.35) {Memory};
\node[box,fill=BitTeal!7,draw=BitTeal!67] (safe) at (7.4,-0.35) {Safety};
\node[box,fill=BitGray,draw=BitSlate!65] (run) at (0,-2.25) {Runtime};
\node[box,fill=BitGray,draw=BitSlate!65] (art) at (3.7,-2.25) {Artifacts};
\node[box,fill=BitGray,draw=BitSlate!65] (cap) at (7.4,-2.25) {Capabilities};
\draw[decorate,decoration={brace,amplitude=5pt},BitSlate] (-0.85,2.22) -- (8.25,2.22) node[midway,above=7pt,figsmall] {hot path};
\draw[decorate,decoration={brace,amplitude=5pt},BitSlate] (-0.85,-2.92) -- (8.25,-2.92) node[midway,below=7pt,figsmall] {cold path};
\end{tikzpicture}
\caption{Classes are distinguished by how they load and how they change, not only by topic.}
\label{fig:grammar-ontology}
\end{figure}



\begin{figure}[H]
\centering
\begin{tikzpicture}[every node/.style={figlabel}, filebox/.style={figbox, align=center, minimum height=.84cm}]
\path[figpanel] (-2.7,-1.95) rectangle (10.25,3.25);
\node[figtitle, anchor=west] at (-2.25,2.8) {TOPOLOGY};
\node[font=\sffamily\bfseries\large, text=BitInk] at (0,2.25) {Single-file monolith};
\node[filebox, fill=BitRose!9, draw=BitRose!62, minimum width=4.2cm, minimum height=2.2cm] (mono) at (0,0.65) {\statefile{AGENTS.md}\\[4pt]\figsmall policy + state + tasks\\history + memory};
\node[align=center, figsmall] at (0,-1.1) {longer context\\mixed lifecycles\\harder review};

\node[font=\sffamily\bfseries\large, text=BitInk] at (6.75,2.25) {CONTINUITY mesh};
\node[filebox, fill=BitBlue!8, draw=BitBlue!65, minimum width=2.15cm] (a) at (4.75,1.35) {\statefile{AGENTS.md}};
\node[filebox, fill=BitTeal!8, draw=BitTeal!70, minimum width=2.15cm] (s) at (6.75,0.7) {\statefile{STATE.md}};
\node[filebox, fill=BitGold!12, draw=BitGold!76!black, minimum width=2.15cm] (t) at (8.75,1.35) {\statefile{TODO.md}};
\node[filebox, fill=BitBlue!8, draw=BitBlue!65, minimum width=2.15cm] (m) at (4.75,0.0) {\statefile{MEMORY.md}};
\node[filebox, fill=BitRose!10, draw=BitRose!62, minimum width=2.15cm] (h) at (8.75,0.0) {\file{completed/}};
\draw[figarrow] (a) -- (s);
\draw[figarrow] (t) -- (s);
\draw[figarrow] (s) -- (h);
\draw[figarrow] (m) -- (s);
\node[align=center, figsmall] at (6.75,-1.1) {split by concern\\loaded by scope\\reviewed by diff};
\end{tikzpicture}
\caption{The mesh preserves separate lifecycles instead of forcing all agent context into one manifest.}
\label{fig:monolith-mesh}
\end{figure}


CONTINUITY organizes repository artifacts into classes with different mutability and load rules. That distinction matters. When current state, historical notes, and reusable knowledge collapse into the same file, the result is usually drift or bloat. In this paper, upper-case filenames such as \statefile{STATE.md} denote the live state layer, while lower-case paths such as \file{completed/} and \file{artifacts/} denote history or evidence stores.


\begin{table}[H]
\centering
\caption{CONTINUITY artifact classes. The key distinction is between what loads by default and what stays cold until needed.}
\label{tab:file-taxonomy}
\small
\rowcolors{2}{BitGray}{white}
\begin{tabularx}{\linewidth}{>{\raggedright\arraybackslash}p{1.55cm}>{\raggedright\arraybackslash}p{3.05cm}>{\raggedright\arraybackslash}p{1.9cm}>{\raggedright\arraybackslash}p{1.75cm}X}
\rowcolor{BitNavy!8}
\textbf{Class} & \textbf{Examples} & \textbf{Mutability} & \textbf{Default load} & \textbf{Function} \\
Boot & \statefile{AGENTS.md}, \statefile{BOOT.md} & rare edits & yes & startup logic, norms, governance \\
State & \statefile{STATE.md}, \statefile{CONTEXT.md} & overwrite & scoped & current project truth \\
Task & \statefile{TODO.md}, \statefile{TASKS.md} & frequent & scoped & executable queue and ownership \\
Memory & \statefile{MEMORY.md}, \statefile{LEARNINGS.md} & curated & selective & durable facts and distilled lessons \\
History & \file{completed/}, \file{snapshots/} & append-only & no & evidence and prior states \\
Safety & \statefile{CHECKPOINT.md}, \statefile{WARNING.md} & stable & yes & approvals and hazard notes \\
Runtime & \file{claims/}, \file{runs/} & structured & per run & activity trace and coordination \\
Artifacts & \file{artifacts/}, \file{reviews/} & append or curated & no & logs, plots, outputs, review records \\
\end{tabularx}
\end{table}


\subsection{Core terms}

\begin{description}[leftmargin=2.6cm,style=nextline]
  \item[Repository-native state] Project state stored inside the repository rather than inside hidden platform memory.
  \item[Current truth] The smallest accurate description of what is true right now.
  \item[State reconciliation] Updating state artifacts after code, tests, or decisions change.
  \item[Scoped loading] Resolving only the manifest, state, and task files relevant to the current path.
  \item[Semantic history] Append-only records organized by meaning: completed work, reviews, evidence, and snapshots.
  \item[Soft constraint] A model-visible instruction that is not technically enforced.
  \item[Hard constraint] A rule enforced by hooks, tests, permissions, schemas, or human approval.
\end{description}

\subsection{Minimal core versus template library}

The paper defends a \term{minimal reference core}, not the largest possible template set. A repository that only needs \statefile{AGENTS.md}, \statefile{STATE.md}, and \statefile{TODO.md} should not be pressured into creating ten more files. The extended topology is valuable as a library of patterns. It is not a proof that every pattern should be loaded.


\begin{figure}[H]
\centering
\begin{tikzpicture}[every node/.style={figlabel}, box/.style={figbox, align=left, minimum width=3.8cm, minimum height=.88cm, inner sep=7pt}]
\path[figpanel] (-2.0,-0.95) rectangle (12.1,3.1);
\node[figtitle, anchor=west] at (-1.55,2.68) {MINIMAL CORE};
\node[font=\sffamily\bfseries\small, text=BitSlate] at (0,2.25) {LIVE STATE};
\node[font=\sffamily\bfseries\small, text=BitSlate] at (5.0,2.25) {MEMORY};
\node[font=\sffamily\bfseries\small, text=BitSlate] at (10.0,2.25) {HISTORY AND EVIDENCE};

\node[box, fill=BitBlue!8, draw=BitBlue!65] at (0,1.4) {\statefile{AGENTS.md}\\\figsmall boot rules};
\node[box, fill=BitTeal!8, draw=BitTeal!70] at (0,0.2) {\statefile{STATE.md}\\\figsmall current truth};
\node[box, fill=BitGold!14, draw=BitGold!76!black] at (0,-1.0) {\statefile{TODO.md}\\\figsmall active queue};
\node[box, fill=BitBlue!8, draw=BitBlue!65] at (5.0,1.4) {\statefile{LEARNINGS.md}\\\figsmall reusable lessons};
\node[box, fill=BitTeal!8, draw=BitTeal!70] at (5.0,0.2) {\statefile{MEMORY.md}\\\figsmall durable facts};
\node[box, fill=BitRose!10, draw=BitRose!62] at (10.0,1.4) {\file{completed/}\\\figsmall finished work};
\node[box, fill=BitGray, draw=BitSlate!65] at (10.0,0.2) {\file{snapshots/}\\\figsmall recovery points};
\node[box, fill=BitGray, draw=BitSlate!65] at (10.0,-1.0) {\file{artifacts/}\\\figsmall evidence};
\end{tikzpicture}
\caption{The minimal core separates live state from archival material. Upper-case filenames denote the live state files.}
\label{fig:minimal-core}
\end{figure}


\subsection{State, history, and knowledge are different objects}

Current state is overwritten. History accumulates. Knowledge is distilled. The operational grammar depends on keeping those three objects separate.


\begin{figure}[H]
\centering
\begin{tikzpicture}[every node/.style={figlabel}, card/.style={figbox, minimum width=3.55cm, minimum height=2.65cm, align=left, inner sep=9pt}]
\path[figpanel] (-1.0,-2.1) rectangle (10.35,2.15);
\node[figtitle, anchor=west] at (-0.55,1.72) {THREE OBJECTS};
\node[card, fill=BitTeal!10, draw=BitTeal!72] (state) at (0,0) {\textbf{STATE}\\[4pt]\figsmall where: root files\\example: \statefile{STATE.md}\\mutability: overwrite\\answers: what is true now};
\node[card, fill=BitRose!10, draw=BitRose!62] (hist) at (4.7,0) {\textbf{HISTORY}\\[4pt]\figsmall where: directories\\example: \file{completed/}\\mutability: append\\answers: what happened};
\node[card, fill=BitBlue!10, draw=BitBlue!62] (lib) at (9.4,0) {\textbf{LIBRARY}\\[4pt]\figsmall where: curated files\\example: \file{memory/}\\mutability: revise\\answers: what is known};
\draw[figarrow] (state.east) -- (hist.west);
\draw[figarrow] (hist.east) -- (lib.west);
\end{tikzpicture}
\caption{Current truth, history, and reusable knowledge should not collapse into one file.}
\label{fig:state-history-knowledge}
\end{figure}

